
\documentclass[conference]{IEEEtran}
\IEEEoverridecommandlockouts
% The preceding line is only needed to identify funding in the first footnote. If that is unneeded, please comment it out.
\usepackage{cite}
\usepackage{amsmath,amssymb,amsfonts}
\usepackage{algorithmic}
\usepackage{graphicx}
\usepackage{textcomp}
\usepackage{xcolor}
\usepackage{listings}
\usepackage{placeins}
\usepackage{booktabs}
\def\BibTeX{{\rm B\kern-.05em{\sc i\kern-.025em b}\kern-.08em
		T\kern-.1667em\lower.7ex\hbox{E}\kern-.125emX}}


\begin{document}
	
	
	\title{Miniproject 4}
	
	\author{\IEEEauthorblockN{Michael Hajostek}
		\IEEEauthorblockA{\textit{School of Electrical Engineering \& Computer Science} \\
			\textit{University of North Dakota}\\
			Grand Forks, ND \\
			michael.hajostek@und.edu}
	}
	
	\maketitle
	
	
	
	\section{Problem 1: Belief Tracking and QMDP}
	
	
	\subsection{1 - Bayesian update by hand}
	
	\subsubsection{A - Written Bayesian update}
	
	
	$b_O(s) = 1/25$ uniform. Drone is at $x_1=(3,3)$ and $o_1=1$.
	
	$b_1(s)=P(o=1|s,x_1)*b_O(s)/Z$
	
	Likelihoods: $P(o=1|s=(3,3), x_1=(3,3))=p_{\text{hit}}=0.9$. 
	Any other cell, $P(o=1|s,x_1=(3,3)) = p_{\text{fp}}=0.1$
	
	Numerators: $(3,3)$ gets $0.9 \times 1/25 = 0.036 + 0.096 = 0.132$
	
	So $b_1((3,3)) = 0.036/0.132 = 9/33 = 3/11 \approx 0.273$ and $b_1(s \neq (3,3)) = 0.004/0.132 = 1/33 \approx 0.030$
	
	Sanity: $3/11 + 24 \times (1/33) = 9/33 + 24/33 = 33/33$.
	
	
	 One ping moves $(3,3)$ from $1/25 = 0.04$ to $\approx 0.273$. Roughly $7\times$ boost.
	
	\subsubsection{B - ten misses, and the asymmetry}
	With each subsequent $o=0$ reading at $(3,3)$, the likelihood ratio for $(3,3)$ vs any other cell is $(1-p_{\text{hit}})/(1-p_{\text{fp}}) = 0.1/0.9 = 1/9$.
	
	So $b((3,3))$ gets multiplied by $1/9$ relative to each other cell on every miss. After 10 misses, the $(3,3)$ posterior collapses by $9^{10} \approx 3.5$ billion relative to the others --- it's effectively ruled out, but the remaining mass spreads across all 24 other cells roughly uniformly (they're symmetric with respect to this measurement sequence).
	
	Explaining the asymmetry: a single hit produces likelihood ratio $p_{\text{hit}}/p_{\text{fp}} = 0.9 / 0.1 = 9$ in favor of the source cell over any one other cell, while a single miss produces ratio $(1-p_{\text{hit}})/(1-p_{\text{fp}}) = 0.1/0.9 = 1/9$ against. Numerically symmetric per cell but asymmetric per belief because:
	
	\begin{itemize}
		\item A hit concentrates probability onto one specific cell against 24 competitors.
		\item A miss redistributes probability across 24 cells that are all equally consistent with the miss.
	\end{itemize}
	
	Thus a hit produces a sharp peak; a miss produces diffuse spreading. The asymmetry is controlled by the likelihood ratios 
	$\frac{p_{\text{hit}}}{p_{\text{fp}}}$ for hits and 
	$\frac{1-p_{\text{hit}}}{1-p_{\text{fp}}}$ for misses. Lower $p_{\text{fp}}$ makes hits much more decisive; the number of alternatives cells then determines why the resulting belief concentration is sharper than the diffuse redistribution cause by misses.	
	
	
	As $p_{\text{fp}} \to 0$, hits become near-certain identifiers (huge concentration) while misses provide weak negative evidence.
	
	As $p_{\text{fp}} \to p_{\text{hit}}$, both observations become uninformative. 
	
	Direction: lower $p_{\text{fp}}$ $\to$ strongest asymmetry favoring fast concentration on hits over slow dispersion on misses.
	
	\subsection{2 - QMDP’s blind spot}
	
	\subsubsection{A - QMDP at $(1,1)$, belief split between $(2,2)$ and $(4,4)$}
	
	QMDP averages $Q^*(s,a)$ under the belief, then arg-maxes. $Q^*$ is shortest-path value to the (assumed-known) target. From $(1,1)$, action ``move toward $(2,2)$'' has high $Q^*$ if $s=(2,2)$ and lowish $Q^*$ if $s=(4,4)$; ``move toward $(4,4)$'' is the reverse. Averaging with weights $0.5/0.5$, the two options come out roughly tied, and QMDP commits to one of them --- say it heads toward $(2,2)$ --- because QMDP assumes uncertainty disappears after one step (it plans as if the next belief is a delta). It will never choose an action whose purpose is to reduce uncertainty.
	
	A belief-aware agent would prefer a probing action:
    choose a first move or route that tests one hypothesis while preserving the ability to redirect if the observation disconfirms it. Under the given sensor model, diagnostic observations occur at the hypothesized cells themselves, not at the midpoint; therefore a belief-aware agent values the future belief update from visiting $(2,2)$ or $(4,4)$, whereas QMDP only averages fully observable shortest-path values and ignores the epistemic value of that test.
	
	\subsubsection{B - Diagnostic pattern in $H(b_t)$ and $a_t$ logs}
	
	The signature of this failure: $H(b_t)$ stays high (belief remains diffuse/multi-modal) while $a_t$ commits to greedy goal-seeking moves toward one mode. Specifically, we'd see entropy plateau or decrease only slowly, while actions are consistently the ``go-to-MAP-cell'' moves rather than actions that would split the hypotheses.
	
	Why this is reliable rather than normal operation: in a healthy POMDP policy, you'd expect either (a) low entropy + committed goal-seeking (belief is concentrated, agent exploits) or (b) high entropy + information-gathering moves (agent hasn't localized yet, so it probes). The pathological combination is high entropy + committed exploitation --- that's QMDP betting on one hypothesis without earning the right to. A correlated diagnostic: track expected information gain of the chosen action vs.\ the max available; QMDP will systematically pick actions with near-zero EIG even when high-EIG alternatives exist.
	
	
	\section*{Acknowledgment}
	
	This material is based upon work supported by the National Science Foundation CISE Graduate
	Fellowships under Grant No. 2313998. Any opinions, findings, and conclusions or
	recommendations expressed in this material are those of the author(s) and do not necessarily
	reflect the views of the National Science Foundation.
	
\end{document}
